% Chapter 10: The Compositional Scale
% DEFINES: ch:compositional-scale; sec:compositional-measures,
%   sec:joint-recovery, sec:irreducibility, sec:defenses, sec:ch10-notes;
%   thm:joint-recovery, rem:assouad-not-lecam, rem:irreducibility, tab:measures
% REFERENCES (backward): thm:orbit-bound, rem:active-vs-marginal, def:acceptance,
%   def:entropy-ratio, prop:fannes, rem:le-cam-marginal, def:totality,
%   thm:composition, rem:granularity
% REFERENCES (forward, part labels): part:constructions, part:frontiers

\chapter{The Compositional Scale}
\label{ch:compositional-scale}

The marginal scale was governed by $\delta$, and a designer who drives
$\delta$ down buys proportional confidentiality for a single cipher value.
This chapter is about the scale where that purchase stops working. When the
untrusted machine evaluates cipher maps, applying the operations it holds,
or observing several maps on the same cipher value, it learns things that no
value of $\delta$ prevents. The headline result is sharp: the latent joint
distribution behind a shared cipher value is recoverable at a rate that no
per-map parameter can move, and the lower bound that makes the rate sharp is
the part's most easily misremembered piece of mathematics. As in
\Cref{ch:marginal-scale}, the results are the Entropy Ratio paper's and are
summarized here with pointers to the full proofs.\footnote{The compositional
theorems and their proofs are \cite{towell2026maxconf}, Theorems 5.1 and 5.2;
this chapter states and motivates, it does not reproduce them.}

\section{Three Ways Operations Leak}
\label{sec:compositional-measures}

The compositional scale is a family of three measures, not one, distinguished
by what the untrusted machine holds (\Cref{tab:measures}). The first is
already in hand. \Cref{thm:orbit-bound} bounded what an \emph{active}
adversary learns by applying its operations to a single token: the residual
entropy $H(X \mid \text{view}) \ge H(X) - \log_2\lvert\orbitF(c)\rvert$,
governed by orbit size, not by $\delta$ (\Cref{rem:active-vs-marginal}). That
is the first member of the family.

The second is the \emph{multi-instance coincidence oracle}. If the untrusted
machine holds $t$ cipher maps for the \emph{same} latent function under
independent secrets, with public acceptance partitions $\{\alpha_i(y)\}$
(\Cref{def:acceptance}), it can look for coincidences, cases where several
instances land in the same value region, and reach accuracy
\[
  \mathrm{accuracy}(t) \;=\; 1 - \tfrac12\sum_{y} \alpha(y)^{t}
\]
in the homogeneous case, growing with the number of instances $t$. The shape
of the acceptance partition is the defense: concentrated partitions, mass on
a few regions, blunt the oracle best, and where shared-function composition
cannot be avoided the remedy is randomized encoding ($K(x) > 1$), not
retuning the partition.\footnote{The coincidence oracle is Theorem 8.2 of
\cite{towell2026cipher}.} The third measure, the shared-variable joint
recovery, is the chapter's headline and the next section.

\begin{table}[t]
\centering
\caption{The three compositional measures. None is controlled by $\delta$.}
\label{tab:measures}
\begin{tabular}{@{}lll@{}}
\toprule
Measure & The untrusted machine holds & Governed by \\
\midrule
Orbit closure & operations to apply to one token & orbit size $\lvert\orbitF(c)\rvert$ \\
Coincidence oracle & $t$ maps for the same $f$ & acceptance-partition shape \\
Joint recovery & maps for $f_1, f_2$ on shared $c$ & latent cardinalities $\lvert Y_1\rvert\lvert Y_2\rvert$ \\
\bottomrule
\end{tabular}
\end{table}

\section{The Joint-Recovery Rate}
\label{sec:joint-recovery}

Suppose the untrusted machine observes pairs $(\fhat_1(c), \fhat_2(c))$,
two cipher maps for two latent functions $f_1 : X \to Y_1$ and
$f_2 : X \to Y_2$, evaluated on the \emph{same} in-domain cipher values $c$.
Two facts hold, and together they are the result.

First, the cipher maps preserve mutual information exactly:
\[
  I\bigl(\fhat_1(C); \fhat_2(C)\bigr) \;=\; I\bigl(f_1(X); f_2(X)\bigr).
\]
The trapdoor hides the values, but it does not destroy the statistical
dependence between them; whatever $f_1(X)$ and $f_2(X)$ shared, their cipher
images still share. Second, an adversary collecting such pairs can estimate
the latent joint distribution on $Y_1 \times Y_2$, and the number of
observations it needs scales as the product of the latent cardinalities.

\begin{theorem}[Joint recovery, summarized]
\label{thm:joint-recovery}
The latent joint distribution on $Y_1 \times Y_2$ is recoverable from shared-%
cipher-value observations to total-variation accuracy $\xi$ at the rate
\[
  N \;=\; \Theta\!\left(\frac{\lvert Y_1\rvert\,\lvert Y_2\rvert}{\xi^{2}}\right),
\]
the upper bound attained by the empirical plug-in estimator and the matching
lower bound by Assouad's lemma. No per-cipher-map parameter, $\delta$
included, changes this rate.
\end{theorem}

The upper bound is unsurprising: the empirical joint distribution of the
observed pairs converges to the latent joint at the usual parametric rate,
and $\lvert Y_1\rvert\lvert Y_2\rvert$ cells take $\lvert Y_1\rvert\lvert
Y_2\rvert/\xi^2$ samples to fill to precision $\xi$. The lower bound is the
subtle half, and its proof technique is the landmine.

\begin{remark}[Assouad, not Le Cam]
\label{rem:assouad-not-lecam}
The matching lower bound is \emph{Assouad's lemma}, not Le Cam's two-point
method. The distinction is not pedantry. Le Cam's two-point method pits two
hypotheses against each other and yields a bound that does not grow with
dimension; it cannot produce the dimension-dependent rate
$\sqrt{\lvert Y_1\rvert\lvert Y_2\rvert / N}$, because that rate comes from
the \emph{many} coordinates of the joint distribution that must be estimated
at once. Assouad's lemma is the right tool: it packs a hypercube of
$2^{m/2}$ hypotheses, one per coordinate of the distribution, and sums the
per-coordinate testing difficulty, yielding exactly the product-of-cardinalities
rate. Le Cam's two-point method does appear in this framework, but at the
marginal scale and as an \emph{upper} bound on a single guesser's accuracy
(\Cref{rem:le-cam-marginal}); using it for this compositional lower bound
would understate the leakage. The two are different theorems with the same
famous name attached, and the compositional one is Assouad.
\end{remark}

To feel the magnitude: recovering the joint of two ten-valued outputs to
total-variation accuracy $\xi = 0.1$ takes on the order of $\lvert
Y_1\rvert\lvert Y_2\rvert/\xi^2 = 100/0.01 = 10^4$ shared-value observations,
a modest number for a server that answers queries in volume.

\section{Why No Parameter Helps}
\label{sec:irreducibility}

The clause that does the work in \Cref{thm:joint-recovery} is the last one:
no per-cipher-map parameter changes the rate. The rate depends only on the
latent cardinalities $\lvert Y_1\rvert$ and $\lvert Y_2\rvert$, on how many
joint cells there are to estimate. It does not depend on $\delta$, on
$\varepsilon$, on $\mu$, on the entropy ratio of \Cref{def:entropy-ratio}, or
on any quantity the marginal scale controls.

\begin{remark}[The two scales do not reduce]
\label{rem:irreducibility}
A small $\delta$ buys a high entropy ratio (\Cref{prop:fannes}) and so caps
marginal leakage; it does \emph{nothing} to the joint-recovery rate. The two
scales measure different things and neither bounds the other: a cipher map
can be marginally near-perfect ($e \approx 1$, $\delta \approx 0$) and still
surrender the full latent joint to an adversary who collects enough shared-%
value pairs. Reducing $\delta$ is therefore necessary but not sufficient for
confidentiality under composition.
\end{remark}

The reason is structural, and it indicts the framework's own virtues.
Totality (\Cref{def:totality}) is what lets the untrusted machine evaluate
any token blindly, and composability (\Cref{thm:composition}) is what lets it
chain evaluations; those two properties are exactly what make a cipher map
useful, and they are exactly what open the compositional channel. The
property that hides a single query is the property that, applied repeatedly
to shared inputs, exposes the joint. There is no setting of the marginal
parameters that closes the channel without also closing off the composition
that motivates the framework.

\section{System-Level Defenses}
\label{sec:defenses}

If the marginal levers do not touch the compositional rate, what does? The
defenses are system-level, acting on the adversary's access rather than on
any single map's parameters. Three reduce the effective leakage. Limiting
the number of shared-cipher-value observations $N$ keeps the adversary below
the sample threshold its accuracy needs, since the rate is in $N$. Raising
the entanglement $p$ of \Cref{rem:granularity}, encoding the outputs of
$f_1$ and $f_2$ jointly rather than separately, hides the very correlation
the joint-recovery attack reconstructs, at the exponential space cost
granularity always carries. Injecting noise slows estimation by the Fisher
factor $\rho^2$ of \Cref{sec:delta-levers}, raising the sample cost without
changing the asymptotic rate. None of these is a marginal-parameter tweak;
all act on the system around the cipher maps. Realizing them concretely,
and weighing their costs, is the work of \Cref{part:constructions}, and the
question of maintaining confidentiality as the query distribution drifts is
among the open problems of \Cref{part:frontiers}.

\section{Notes and Provenance}
\label{sec:ch10-notes}

The three compositional measures are the framework's, collected from the
Algebraic Cipher Types paper (orbit closure), the Cipher Maps paper
(coincidence oracle), and the Entropy Ratio paper (joint recovery); the
headline joint-recovery rate and its bounds are the last paper's Theorems 5.1
and 5.2, summarized here with the proofs left to \cite{towell2026maxconf}.
The chapter carries two more of the book's landmines. The compositional lower
bound is \emph{Assouad's lemma}, over a $2^{m/2}$ hypercube packing, not Le
Cam's two-point method, and the two roles of the two-point method, the
marginal upper bound of \Cref{rem:le-cam-marginal} and this compositional
lower bound, must be kept apart. And the residual-entropy and joint-recovery
quantities are unnormalized leakage measures, not the normalized entropy
ratio of the marginal scale, and the two should never be equated. With this
chapter the two-scale confidentiality theory the book has been building since
\Cref{sec:measured-not-negligible} is complete: a marginal scale that
$\delta$ governs, a compositional scale that it does not, and a clear account
of which defenses act where.
